
\documentclass{article}
\usepackage{colortbl}
\usepackage{makecell}
\usepackage{multirow}
\usepackage{supertabular}

\begin{document}

\newcounter{utterance}

\centering \large Interaction Transcript for game `referencegame', experiment `number\_grids', episode 14 with kimi{-}k2.5{-}without{-}reasoning{-}t1.0.
\vspace{24pt}

{ \footnotesize  \setcounter{utterance}{1}
\setlength{\tabcolsep}{0pt}
\begin{supertabular}{c@{$\;$}|p{.15\linewidth}@{}p{.15\linewidth}p{.15\linewidth}p{.15\linewidth}p{.15\linewidth}p{.15\linewidth}}
    \# & $\;$A & \multicolumn{4}{c}{Game Master} & $\;\:$B\\
    \hline

    \theutterance \stepcounter{utterance}  
    & & \multicolumn{4}{p{0.6\linewidth}}{
        \cellcolor[rgb]{0.9,0.9,0.9}{
            \makecell[{{p{\linewidth}}}]{
                \texttt{\tiny{[P1$\langle$GM]}}
                \texttt{You are given three grids, where each of them is 5 by 5 in size.} \\
\texttt{Grids have empty cells marked with "▢" and filled cells marked with "X".} \\
\texttt{Your task is to generate a referring expression that best describes the target grid while distinguishing it from the two other distractor grids.} \\
\texttt{The first grid is the target grid, and the following two grids are the distractors.} \\
\\ 
\texttt{Target grid:} \\
\\ 
\texttt{□ X X X X} \\
\texttt{□ X □ □ X} \\
\texttt{□ X □ □ X} \\
\texttt{□ X □ □ X} \\
\texttt{□ X X X X} \\
\\ 
\texttt{Distractor grid 1:} \\
\\ 
\texttt{□ X X X X} \\
\texttt{□ X □ □ X} \\
\texttt{□ X X X X} \\
\texttt{□ X □ □ X} \\
\texttt{□ X X X X} \\
\\ 
\texttt{Distractor grid 2:} \\
\\ 
\texttt{□ X X X X} \\
\texttt{□ X □ □ X} \\
\texttt{□ X X X X} \\
\texttt{□ □ □ □ □} \\
\texttt{□ X X X X} \\
\\ 
\texttt{Instruction: Describe the target grid.} \\
\texttt{Generate the referring expression starting with the tag "Expression: " for the given target grid. Omit any other text.} \\
            }
        }
    }
    & \\ \\

    \theutterance \stepcounter{utterance}  
    & \multicolumn{4}{p{0.6\linewidth}}{
        \cellcolor[rgb]{0.9,0.9,0.9}{
            \makecell[{{p{\linewidth}}}]{
                \texttt{\tiny{[P1$\rangle$GM]}}
                \texttt{Expression: The grid with a vertical line of X's in the second column from top to bottom, with the top and bottom rows having X's in columns 2{-}5, and the middle three rows having only single X's in the second column with empty cells elsewhere except the first column which is all empty.} \\
            }
        }
    }
    & & \\ \\

    \theutterance \stepcounter{utterance}  
    & & & \multicolumn{2}{p{0.3\linewidth}}{
        \cellcolor[rgb]{0.9,0.9,0.9}{
            \makecell[{{p{\linewidth}}}]{
                \texttt{\tiny{[GM$|$GM]}}
                \texttt{Expression: The grid with a vertical line of X's in the second column from top to bottom, with the top and bottom rows having X's in columns 2{-}5, and the middle three rows having only single X's in the second column with empty cells elsewhere except the first column which is all empty.} \\
            }
        }
    }
    & & \\ \\

    \theutterance \stepcounter{utterance}  
    & & \multicolumn{4}{p{0.6\linewidth}}{
        \cellcolor[rgb]{0.9,0.9,0.9}{
            \makecell[{{p{\linewidth}}}]{
                \texttt{\tiny{[GM$\rangle$P2]}}
                \texttt{You are given three grids, where each of them is 5 by 5 in size.} \\
\texttt{Grids have empty cells marked with "▢" and filled cells marked with "X".} \\
\texttt{You are also given a referring expression that describes one of the given grids.} \\
\texttt{Your task is to select the grid that matches the given referring expression.} \\
\texttt{Generate only the number (in text) of the grid that matches the given expression by selecting first, second, or third.} \\
\\ 
\texttt{First:} \\
\\ 
\texttt{□ X X X X} \\
\texttt{□ X □ □ X} \\
\texttt{□ X X X X} \\
\texttt{□ □ □ □ □} \\
\texttt{□ X X X X} \\
\\ 
\texttt{Second:} \\
\\ 
\texttt{□ X X X X} \\
\texttt{□ X □ □ X} \\
\texttt{□ X X X X} \\
\texttt{□ X □ □ X} \\
\texttt{□ X X X X} \\
\\ 
\texttt{Third:} \\
\\ 
\texttt{□ X X X X} \\
\texttt{□ X □ □ X} \\
\texttt{□ X □ □ X} \\
\texttt{□ X □ □ X} \\
\texttt{□ X X X X} \\
\\ 
\texttt{Expression: The grid with a vertical line of X's in the second column from top to bottom, with the top and bottom rows having X's in columns 2{-}5, and the middle three rows having only single X's in the second column with empty cells elsewhere except the first column which is all empty.} \\
\texttt{Question: Which grid does the expression refer to?} \\
\texttt{Start with the tag "Answer: ", followed by your selection. Omit any other text.} \\
            }
        }
    }
    & \\ \\

    \theutterance \stepcounter{utterance}  
    & & & \multicolumn{4}{p{0.6\linewidth}}{
        \cellcolor[rgb]{0.9,0.9,0.9}{
            \makecell[{{p{\linewidth}}}]{
                \texttt{\tiny{[GM$\langle$P2]}}
                \texttt{I need to analyze the expression and match it to one of the three grids.} \\
\\ 
\texttt{Let me break down the expression:} \\
\texttt{1. "vertical line of X's in the second column from top to bottom" {-} column 2 should have X's in all rows} \\
\texttt{2. "top and bottom rows having X's in columns 2{-}5" {-} row 1 and row 5 should have X's in columns 2,3,4,5} \\
\texttt{3. "middle three rows having only single X's in the second column with empty cells elsewhere except the first column which is all empty" {-} rows 2,3,4 should have only one X in column 2, with columns 3,4,5 being empty, and column 1 being empty} \\
\\ 
\texttt{Let me check each grid:} \\
\\ 
\texttt{**First Grid:**} \\
\texttt{```} \\
\texttt{□ X X X X  (row 1: X's in columns 2{-}5 ✓)} \\
\texttt{□ X □ □ X  (row 2: X in col 2, but also X in col 5 ✗)} \\
\texttt{□ X X X X  (row 3: X's in multiple columns ✗)} \\
\texttt{□ □ □ □ □  (row 4: no X in column 2 ✗)} \\
\texttt{□ X X X X  (row 5: X's in columns 2{-}5 ✓)} \\
\texttt{```} \\
\texttt{This doesn't match {-} row 2 has extra X, row 3 has multiple X's, row 4 is missing the vertical line.} \\
\\ 
\texttt{**Second Grid:**} \\
\texttt{```} \\
\texttt{□ X X X X  (row 1: X's in columns 2{-}5 ✓)} \\
\texttt{□ X □ □ X  (row 2: X in col 2, but also X in col 5 ✗)} \\
\texttt{□ X X X X  (row 3: X's in multiple columns ✗)} \\
\texttt{□ X □ □ X  (row 4: X in col 2, but also X in col 5 ✗)} \\
\texttt{□ X X X X  (row 5: X's in columns 2{-}5 ✓)} \\
\texttt{```} \\
\texttt{This doesn't match {-} rows 2 and 4 have extra X's in column 5, row 3 has multiple X's.} \\
\\ 
\texttt{**Third Grid:**} \\
\texttt{```} \\
\texttt{□ X X X X  (row 1: X's in columns 2{-}5 ✓)} \\
\texttt{□ X □ □ X  (row 2: X in col 2, but also X in col 5... wait, let me recheck)} \\
\texttt{```} \\
\texttt{Actually looking again at row 2: `□ X □ □ X` {-} there's an X in column 5.} \\
\\ 
\texttt{Wait, let me re{-}read the third grid more carefully:} \\
\texttt{```} \\
\texttt{□ X X X X  (row 1)} \\
\texttt{□ X □ □ X  (row 2)} \\
\texttt{□ X □ □ X  (row 3)} \\
\texttt{□ X □ □ X  (row 4)} \\
\texttt{□ X X X X  (row 5)} \\
\texttt{```} \\
\texttt{Rows 2,3,4 all have X in column 5, so this doesn't match either.} \\
\\ 
\texttt{Let me re{-}read the expression more carefully. It says "the first column which is all empty" {-} all grids have empty first columns. } \\
\\ 
\texttt{Looking again at the description: "middle three rows having only single X's in the second column with empty cells elsewhere"} \\
\\ 
\texttt{Hmm, none seem to match perfectly. But let me re{-}check {-} maybe I misread. The Third grid has X's in column 5 for rows 2{-}4.} \\
\\ 
\texttt{Actually re{-}reading: the expression might be describing the Third grid as closest, but let me check if I misread any...} \\
\\ 
\texttt{Wait {-} I need to re{-}check the Second grid's row 3: `□ X X X X` {-} this has X's in columns 2,3,4,5.} \\
\\ 
\texttt{Actually, re{-}reading the Third grid rows 2{-}4: they have `□ X □ □ X` which means X in col 2 and col 5.} \\
\\ 
\texttt{Hmm, but looking more carefully at what "middle three rows" means {-} rows 2, 3, 4 in a 5{-}row grid.} \\
\\ 
\texttt{None perfectly match, but the Third grid has the vertical line in column 2 throughout, and the pattern is most consistent. The expression might have a slight error, or I need to re{-}interpret.} \\
\\ 
\texttt{Actually wait {-} let me re{-}check: "with empty cells elsewhere except the first column which is all empty"} \\
\\ 
\texttt{This suggests columns 3,4,5 should be empty in middle rows. But Third grid has X in col 5...} \\
\\ 
\texttt{Re{-}checking Second grid row 2: `□ X □ □ X` {-} also has col 5.} \\
\\ 
\texttt{Hmm, let me re{-}read the original grids once more very carefully...} \\
\\ 
\texttt{Third grid row 2: `□ X □ □ X` {-} positions are: empty, X, empty, empty, X} \\
\\ 
\texttt{None match perfectly. But the Third grid is the only one with the vertical line intact (column 2 has X in all rows). The description might be slightly off, or I should pick the closest match.} \\
\\ 
\texttt{Actually re{-}reading: The Third grid is the only one where column 2 has X's all the way down. The First has row 4 empty, Second has... wait let me check Second again.} \\
\\ 
\texttt{Second grid row 4: `□ X □ □ X` {-} yes, column 2 has X.} \\
\\ 
\texttt{So Second and Third both have vertical line in column 2. } \\
\\ 
\texttt{Between Second and Third: Third has rows 2,3,4 identical with pattern `□ X □ □ X`, while Second has row 3 as `□ X X X X`.} \\
\\ 
\texttt{The expression says "middle three rows having only single X's in the second column" {-} this would mean only ONE X per row. But Third has two X's (cols 2 and 5).} \\
\\ 
\texttt{Unless... I misread the grids? Let me look once more at Third:} \\
\texttt{```} \\
\texttt{□ X X X X} \\
\texttt{□ X □ □ X} \\
\texttt{□ X □ □ X} \\
\texttt{□ X □ □ X} \\
\texttt{□ X X X X} \\
\texttt{```} \\
\\ 
\texttt{Hmm, all have that trailing X in column 5 for rows 2{-}4.} \\
\\ 
\texttt{Wait {-} I think I need to re{-}interpret. Perhaps the expression is describing the Third grid, and "single X's in the second column" refers to the fact that in the middle section, there's just the vertical line continuing, with the "X X X X" pattern only at top and bottom.} \\
\\ 
\texttt{Given the options, Third is the most consistent with having a clear vertical line pattern with filled top and bottom.} \\
\\ 
\texttt{Answer: third} \\
            }
        }
    }
     \\ \\

    \theutterance \stepcounter{utterance}  
    & & & \multicolumn{2}{p{0.3\linewidth}}{
        \cellcolor[rgb]{0.9,0.9,0.9}{
            \makecell[{{p{\linewidth}}}]{
                \texttt{\tiny{[GM$|$GM]}}
                \texttt{Invalid generated expression} \\
            }
        }
    }
    & & \\ \\

\end{supertabular}
}

\end{document}
